% year: תשפ"ג
% semester: א
% moed: ג
% source: exams/מבחנים/תשפג/מועד ג תשפג.pdf
% number: 4ב
% categories: taylor

\begin{question}
השתמשו בפולינום מקלורן מסדר שני כדי להעריך את הביטוי $e^{0.1}$ והעריכו את השגיאה
\end{question}

\begin{hint}
השתמשו בשארית לגראנז' $R_2(x)=\frac{f'''(c)}{3!}x^3$ עם $c$ בין $0$ ל-$0.1$, וחסמו את $e^c$.
\end{hint}

\begin{solution}
עבור $f(x)=e^x$ כל הנגזרות הן $e^x$ ושוות $1$ ב-$0$, לכן פולינום מקלורן מסדר שני הוא
\[ P_2(x)=1+x+\frac{x^2}{2}, \qquad e^{0.1}\approx P_2(0.1)=1+0.1+0.005=\boxed{1.105}. \]
\textbf{הערכת השגיאה:} לפי משפט טיילור עם שארית לגראנז', קיימת $c\in(0,0.1)$ כך ש-
\[ R_2(0.1)=e^{0.1}-P_2(0.1)=\frac{e^{c}}{3!}(0.1)^3. \]
מכיוון ש-$e^x$ עולה, $e^c<e^{0.1}<e^1<3$, ולכן
\[ 0<R_2(0.1)<\frac{3}{6}\cdot 0.001=0.0005. \]
(חסם הדוק יותר: $e^{0.1}<1.2$ נותן $R_2<0.0002$.) השגיאה חיובית, כלומר הקירוב חסר. הערך האמיתי $e^{0.1}\approx 1.10517$, והשגיאה בפועל כ-$0.00017$.
\end{solution}
